% !TEX program = xelatex
% Compile with XeLaTeX. Requires the standard ctex/Fandol packages in TeX Live.
% 中文之后为英文；独立公式共用。Chinese text is followed by English.
\documentclass[11pt,a4paper]{article}

% --- Page margins and paragraphs (based on the supplied template) ---
\usepackage[margin=1in]{geometry}
\usepackage{parskip}
\usepackage[UTF8,scheme=plain,fontset=fandol]{ctex}
\usepackage{fontspec}
\setmainfont{Latin Modern Roman}

% --- Basic mathematics and typesetting ---
\usepackage{amsmath,amssymb,amsthm,mathtools}
\usepackage{xcolor}
\usepackage{longtable,booktabs,array,calc}
\usepackage{enumitem}
\usepackage{fancyhdr}
\usepackage{titlesec}
\usepackage{needspace}
\usepackage[hidelinks,unicode]{hyperref}
\definecolor{TranslationBlue}{HTML}{36546C}
\definecolor{NoteGray}{HTML}{505963}

% --- Definition and theorem environments, retained for further editing ---
\theoremstyle{definition}
\newtheorem{definition}{定义 / Definition}
\newtheorem{example}{例 / Example}
\theoremstyle{plain}
\newtheorem{theorem}{定理 / Theorem}
\newtheorem{proposition}{命题 / Proposition}
\newtheorem{lemma}{引理 / Lemma}

% --- Bilingual paragraph styles ---
\newenvironment{englishtranslation}
  {\par\begingroup\small\color{TranslationBlue}\setlength{\parskip}{3pt}}
  {\par\endgroup}
\newenvironment{editorialnote}
  {\par\begingroup\small\color{NoteGray}\setlength{\parskip}{4pt}}
  {\par\endgroup\medskip}
\newenvironment{bilingualpair}
  {\par\noindent\begin{minipage}{\linewidth}\setlength{\parskip}{4pt}\interdisplaylinepenalty=10000}
  {\end{minipage}\par\addvspace{4pt}}
\setlength{\parskip}{4pt plus 1pt}
\setlength{\parindent}{0pt}
\linespread{1.06}
\setlength{\emergencystretch}{2em}
\setlength{\abovedisplayskip}{8pt plus 2pt minus 2pt}
\setlength{\belowdisplayskip}{8pt plus 2pt minus 2pt}
\setlength{\jot}{4pt}
\allowdisplaybreaks[1]
\raggedbottom
\setcounter{secnumdepth}{0}
\setlist{topsep=4pt,itemsep=3pt,parsep=2pt,leftmargin=1.6em}
\providecommand{\tightlist}{\setlength{\itemsep}{0pt}\setlength{\parskip}{0pt}}
\providecommand{\passthrough}[1]{#1}

\titleformat{\section}{\large\bfseries}{}{0pt}{}
\titleformat{\subsection}{\normalsize\bfseries}{}{0pt}{}
\titleformat{\subsubsection}{\normalsize\bfseries}{}{0pt}{}
\titlespacing*{\section}{0pt}{16pt plus 3pt}{8pt}
\titlespacing*{\subsection}{0pt}{12pt plus 2pt}{6pt}
\titlespacing*{\subsubsection}{0pt}{10pt plus 2pt}{5pt}
\pagestyle{fancy}
\fancyhf{}
\fancyhead[L]{\small 数学物理方法 / Mathematical Methods for Physics}
\fancyhead[R]{\small Lecture 15}
\fancyfoot[C]{\small\thepage}
\setlength{\headheight}{15pt}
\renewcommand{\headrulewidth}{0.3pt}
\hypersetup{pdftitle={Mathematical Methods for Physics - Lecture 15}}

\begin{document}
\begin{center}
  {\LARGE\bfseries 数学物理方法 · 第 15 讲}\\[0.5em]
  {\Large\bfseries Mathematical Methods for Physics}\\[0.2em]
  {\large Lecture 15}\\[0.7em]
  {\normalsize 卷积、傅里叶变换的应用与拉普拉斯变换}\\[0.2em]
  {\small\color{TranslationBlue} Convolution, Applications of Fourier Transforms, and the Laplace Transform}
\end{center}
\vspace{0.5em}

\Needspace{12\baselineskip}

\section{\texorpdfstring{一、卷积与傅里叶变换\\[2pt]{\normalfont\normalsize\color{TranslationBlue} Convolution
and Fourier
Transforms}}{一、卷积与傅里叶变换Convolution and Fourier Transforms}}

\Needspace{8\baselineskip}

\subsection{\texorpdfstring{1.
卷积的定义及运算性质\\[2pt]{\normalfont\small\color{TranslationBlue} Definition
and Algebraic Properties of
Convolution}}{1. 卷积的定义及运算性质Definition and Algebraic Properties of Convolution}}

\begin{bilingualpair}

\emph{接第 14 讲的``性质 8：卷积性''。}

\nopagebreak[4]\begin{englishtranslation}

\emph{continuation of
``Property 8: Convolution'' in Lecture 14.}

\end{englishtranslation}

\end{bilingualpair}

\begin{bilingualpair}

\textbf{定义（卷积）}

\nopagebreak[4]\begin{englishtranslation}

\textbf{Definition (convolution)}

\end{englishtranslation}

\end{bilingualpair}

\begin{bilingualpair}

定义在 \((-\infty,+\infty)\) 上的两个函数
\(f_1(t)\)、\(f_2(t)\)，其卷积定义为

\nopagebreak[4]\begin{englishtranslation}

For two functions \(f_1(t)\) and \(f_2(t)\) defined on
\((-\infty,+\infty)\), their convolution is defined by

\end{englishtranslation}

\[
(f_1*f_2)(t)
=\int_{-\infty}^{+\infty}
f_1(\tau)f_2(t-\tau)\,\mathrm d\tau.
\]

\end{bilingualpair}

\begin{bilingualpair}

\textbf{运算性质}

\nopagebreak[4]\begin{englishtranslation}

\textbf{Algebraic properties}

\end{englishtranslation}

\[
f_1*f_2=f_2*f_1,
\]

\end{bilingualpair}

\[
f_1*(f_2*f_3)=(f_1*f_2)*f_3,
\]

\[
f_1*(f_2+f_3)=f_1*f_2+f_1*f_3.
\]

\begin{bilingualpair}
Notes:
\par
\vspace*{3cm}
\end{bilingualpair}

\Needspace{8\baselineskip}

\subsection{\texorpdfstring{2.
卷积定理\\[2pt]{\normalfont\small\color{TranslationBlue} The Convolution
Theorem}}{2. 卷积定理The Convolution Theorem}}

\begin{bilingualpair}

设

\nopagebreak[4]\begin{englishtranslation}

Let

\end{englishtranslation}

\[
\mathcal F[f_1(t)]=F_1(\omega),
\qquad
\mathcal F[f_2(t)]=F_2(\omega).
\]

\end{bilingualpair}

\begin{bilingualpair}

\textbf{（1）卷积的傅里叶变换等于变换的乘积}

\nopagebreak[4]\begin{englishtranslation}

\textbf{(1) The Fourier transform of a convolution is the product of the
transforms}

\end{englishtranslation}

\[
\mathcal F[(f_1*f_2)(t)]
=\mathcal F[f_1]\cdot\mathcal F[f_2]
=F_1(\omega)F_2(\omega).
\]

\end{bilingualpair}

\begin{bilingualpair}

即

\nopagebreak[4]\begin{englishtranslation}

That is,

\end{englishtranslation}

\[
f_1*f_2
\ \xleftrightarrow{\mathcal F}\
F_1F_2.
\]

\end{bilingualpair}

\begin{bilingualpair}

\textbf{（2）乘积的傅里叶变换}

\nopagebreak[4]\begin{englishtranslation}

\textbf{(2) Fourier transform of a product}

\end{englishtranslation}

\[
\mathcal F[f_1f_2]
=\frac1{2\pi}\bigl(\mathcal F[f_1]*\mathcal F[f_2]\bigr)
=\frac1{2\pi}(F_1*F_2).
\]

\end{bilingualpair}

\begin{bilingualpair}

\textbf{证明（1）}

\nopagebreak[4]\begin{englishtranslation}

\textbf{Proof of (1)}

\end{englishtranslation}

\[
\begin{aligned}
\mathcal F[(f_1*f_2)(t)]
&=\int_{-\infty}^{+\infty}
(f_1*f_2)(t)e^{-i\omega t}\,\mathrm dt\\
&=\int_{-\infty}^{+\infty}
\left[
\int_{-\infty}^{+\infty}
f_1(\tau)f_2(t-\tau)\,\mathrm d\tau
\right]e^{-i\omega t}\,\mathrm dt\\
&=\int_{-\infty}^{+\infty}
f_1(\tau)
\left[
\int_{-\infty}^{+\infty}
f_2(t-\tau)e^{-i\omega t}\,\mathrm dt
\right]\mathrm d\tau.
\end{aligned}
\]

\end{bilingualpair}

\begin{bilingualpair}

利用平移性，

\nopagebreak[4]\begin{englishtranslation}

By the shifting property,

\end{englishtranslation}

\[
\int_{-\infty}^{+\infty}
f_2(t-\tau)e^{-i\omega t}\,\mathrm dt
=e^{-i\tau\omega}F_2(\omega).
\]

\end{bilingualpair}

\begin{bilingualpair}

因此

\nopagebreak[4]\begin{englishtranslation}

Therefore,

\end{englishtranslation}

\[
\begin{aligned}
\mathcal F[(f_1*f_2)(t)]
&=\int_{-\infty}^{+\infty}
f_1(\tau)e^{-i\tau\omega}F_2(\omega)\,\mathrm d\tau\\
&=F_2(\omega)
\int_{-\infty}^{+\infty}
f_1(\tau)e^{-i\tau\omega}\,\mathrm d\tau\\
&=F_2(\omega)F_1(\omega).
\end{aligned}
\]

\end{bilingualpair}

\begin{bilingualpair}
Notes:
\par
\vspace*{3cm}
\end{bilingualpair}

\Needspace{8\baselineskip}

\subsection{\texorpdfstring{3.
例：用卷积证明积分性质\\[2pt]{\normalfont\small\color{TranslationBlue} Example:
Proving the Integration Property by
Convolution}}{3. 例：用卷积证明积分性质Example: Proving the Integration Property by Convolution}}

\begin{bilingualpair}

若

\nopagebreak[4]\begin{englishtranslation}

If

\end{englishtranslation}

\[
\int_{-\infty}^{+\infty}f(t)\,\mathrm dt\ne0,
\]

\end{bilingualpair}

\begin{bilingualpair}

则

\nopagebreak[4]\begin{englishtranslation}

then

\end{englishtranslation}

\[
\mathcal F\!\left[
\int_{-\infty}^{t}f(\tau)\,\mathrm d\tau
\right]
=\frac{F(\omega)}{i\omega}
+\pi\delta(\omega)F(0).
\]

\end{bilingualpair}

\begin{bilingualpair}

\textbf{证明}

\nopagebreak[4]\begin{englishtranslation}

\textbf{Proof}

\end{englishtranslation}

\end{bilingualpair}

\begin{bilingualpair}

先说明

\nopagebreak[4]\begin{englishtranslation}

First, we show that

\end{englishtranslation}

\[
\int_{-\infty}^{t}f(\tau)\,\mathrm d\tau
=(f*u)(t),
\]

\end{bilingualpair}

\begin{bilingualpair}

其中

\nopagebreak[4]\begin{englishtranslation}

where

\end{englishtranslation}

\[
u(t)=
\begin{cases}
1,&t>0,\\
0,&t<0.
\end{cases}
\]

\end{bilingualpair}

\begin{bilingualpair}

事实上，

\nopagebreak[4]\begin{englishtranslation}

Indeed,

\end{englishtranslation}

\[
(f*u)(t)
=\int_{-\infty}^{+\infty}
f(\tau)u(t-\tau)\,\mathrm d\tau.
\]

\end{bilingualpair}

\begin{bilingualpair}

令 \(t-\tau=\xi\)，则

\nopagebreak[4]\begin{englishtranslation}

Let \(t-\tau=\xi\). Then

\end{englishtranslation}

\[
\begin{aligned}
(f*u)(t)
&=\int_{+\infty}^{-\infty}
f(t-\xi)u(\xi)(-\mathrm d\xi)\\
&=\int_{-\infty}^{+\infty}
f(t-\xi)u(\xi)\,\mathrm d\xi\\
&=\int_{-\infty}^{0}0\,\mathrm d\xi
+\int_0^{+\infty}f(t-\xi)\,\mathrm d\xi.
\end{aligned}
\]

\end{bilingualpair}

\begin{bilingualpair}

再令 \(t-\xi=\tau\)，得

\nopagebreak[4]\begin{englishtranslation}

Now let \(t-\xi=\tau\). We obtain

\end{englishtranslation}

\[
(f*u)(t)
=\int_t^{-\infty}f(\tau)(-\mathrm d\tau)
=\int_{-\infty}^{t}f(\tau)\,\mathrm d\tau.
\]

\end{bilingualpair}

\begin{bilingualpair}

由卷积定理，

\nopagebreak[4]\begin{englishtranslation}

By the convolution theorem,

\end{englishtranslation}

\[
\begin{aligned}
\mathcal F\!\left[
\int_{-\infty}^{t}f(\tau)\,\mathrm d\tau
\right]
&=\mathcal F[(f*u)(t)]\\
&=\mathcal F[f]\cdot\mathcal F[u]\\
&=F(\omega)
\left[\frac1{i\omega}+\pi\delta(\omega)\right]\\
&=\frac{F(\omega)}{i\omega}
+F(\omega)\pi\delta(\omega)\\
&=\frac{F(\omega)}{i\omega}
+F(0)\pi\delta(\omega).
\end{aligned}
\]

\end{bilingualpair}

\begin{bilingualpair}
Notes:
\par
\vspace*{3cm}
\end{bilingualpair}

\Needspace{12\baselineskip}

\section{\texorpdfstring{二、傅里叶变换的应用\\[2pt]{\normalfont\normalsize\color{TranslationBlue} Applications
of Fourier
Transforms}}{二、傅里叶变换的应用Applications of Fourier Transforms}}

\begin{bilingualpair}

\emph{教材 P.200。}

\nopagebreak[4]\begin{englishtranslation}

\emph{textbook, p.~200.}

\end{englishtranslation}

\end{bilingualpair}

\begin{bilingualpair}

原稿列出两类应用：

\nopagebreak[4]\begin{englishtranslation}

The original notes list two types of application.

\end{englishtranslation}

\end{bilingualpair}

\begin{itemize}
\item
  \begin{bilingualpair}

  求解微分方程。

  \nopagebreak[4]\begin{englishtranslation}

  Solving differential equations.

  \end{englishtranslation}

  \end{bilingualpair}
\item
  \begin{bilingualpair}

  求解积分方程。

  \nopagebreak[4]\begin{englishtranslation}

  Solving integral equations.

  \end{englishtranslation}

  \end{bilingualpair}
\end{itemize}

\Needspace{8\baselineskip}

\subsection{\texorpdfstring{1. 例
7.20：求解微分方程\\[2pt]{\normalfont\small\color{TranslationBlue} Example
7.20: Solving a Differential
Equation}}{1. 例 7.20：求解微分方程Example 7.20: Solving a Differential Equation}}

\begin{bilingualpair}

求解

\nopagebreak[4]\begin{englishtranslation}

Solve

\end{englishtranslation}

\[
x''(t)-x(t)=-f(t),
\qquad -\infty<t<+\infty.
\]

\end{bilingualpair}

\begin{bilingualpair}

原稿蓝字补充：为使用微分性，假设

\nopagebreak[4]\begin{englishtranslation}

An annotation in blue in the original introduces the following
assumptions for applying the differentiation property.

\end{englishtranslation}

\[
x(\pm\infty)=0,
\qquad
x'(\pm\infty)=0.
\]

\end{bilingualpair}

\begin{bilingualpair}

\textbf{解}

\nopagebreak[4]\begin{englishtranslation}

\textbf{Solution}

\end{englishtranslation}

\end{bilingualpair}

\begin{bilingualpair}

\textbf{第一步，取傅里叶变换。}

\nopagebreak[4]\begin{englishtranslation}

\textbf{Step 1. Take the Fourier transform.}

\end{englishtranslation}

\end{bilingualpair}

\begin{bilingualpair}

设

\nopagebreak[4]\begin{englishtranslation}

Let

\end{englishtranslation}

\[
\mathcal F[x(t)]=\overline X(\omega),
\qquad
\mathcal F[f(t)]=F(\omega).
\]

\end{bilingualpair}

\begin{bilingualpair}

对方程两边取变换，

\nopagebreak[4]\begin{englishtranslation}

Taking the Fourier transform of both sides gives

\end{englishtranslation}

\[
\mathcal F[x''(t)-x(t)]=-\mathcal F[f(t)].
\]

\end{bilingualpair}

\begin{bilingualpair}

由微分性，

\nopagebreak[4]\begin{englishtranslation}

By the differentiation property,

\end{englishtranslation}

\[
(i\omega)^2\overline X(\omega)-\overline X(\omega)
=-F(\omega),
\]

\end{bilingualpair}

\begin{bilingualpair}

所以

\nopagebreak[4]\begin{englishtranslation}

Therefore,

\end{englishtranslation}

\[
\overline X(\omega)=\frac{F(\omega)}{\omega^2+1}.
\]

\end{bilingualpair}

\begin{bilingualpair}

\textbf{第二步，取傅里叶反变换。}

\nopagebreak[4]\begin{englishtranslation}

\textbf{Step 2. Take the inverse Fourier transform.}

\end{englishtranslation}

\[
\begin{aligned}
x(t)
&=\mathcal F^{-1}[\overline X(\omega)]\\
&=\mathcal F^{-1}\!\left[
F(\omega)\frac1{\omega^2+1}
\right]\\
&=\mathcal F^{-1}[F(\omega)]
*\mathcal F^{-1}\!\left[\frac1{\omega^2+1}\right]\\
&=f(t)*\left(\frac12e^{-|t|}\right)\\
&=\frac12\int_{-\infty}^{+\infty}
f(\tau)e^{-|t-\tau|}\,\mathrm d\tau.
\end{aligned}
\]

\end{bilingualpair}

\begin{bilingualpair}

这就是原稿所求的解。

\nopagebreak[4]\begin{englishtranslation}

This is the solution obtained in the original notes.

\end{englishtranslation}

\end{bilingualpair}

\begin{bilingualpair}

其中使用了变换对

\nopagebreak[4]\begin{englishtranslation}

The transform pair used here is

\end{englishtranslation}

\[
e^{-\beta|t|}
\ \xleftrightarrow{\mathcal F}\
\frac{2\beta}{\omega^2+\beta^2},
\]

\end{bilingualpair}

\begin{bilingualpair}

从而

\nopagebreak[4]\begin{englishtranslation}

It follows that

\end{englishtranslation}

\[
\mathcal F^{-1}\!\left[\frac1{\omega^2+\beta^2}\right]
=\frac1{2\beta}e^{-\beta|t|},
\qquad \beta>0.
\]

\end{bilingualpair}

\begin{editorialnote}

\begin{bilingualpair}

\textbf{整理注}：\(\beta>0\)
为此衰减指数变换对的条件。

\nopagebreak[4]\begin{englishtranslation}

\textbf{Editorial note.} The condition \(\beta>0\) is required for this
decaying-exponential transform pair.

\end{englishtranslation}

\end{bilingualpair}

\end{editorialnote}

\begin{bilingualpair}
Notes:
\par
\vspace*{3cm}
\end{bilingualpair}

\Needspace{8\baselineskip}

\subsection{\texorpdfstring{2. 例
7.22：求解卷积型积分方程\\[2pt]{\normalfont\small\color{TranslationBlue} Example
7.22: Solving a Convolution Integral
Equation}}{2. 例 7.22：求解卷积型积分方程Example 7.22: Solving a Convolution Integral Equation}}

\begin{bilingualpair}

求解积分方程

\nopagebreak[4]\begin{englishtranslation}

Solve the integral equation

\end{englishtranslation}

\[
\int_{-\infty}^{+\infty}
e^{-|t-\xi|}x(\xi)\,\mathrm d\xi
=f(t),
\]

\end{bilingualpair}

\begin{bilingualpair}

其中 \(f(t)\) 已知。

\nopagebreak[4]\begin{englishtranslation}

where \(f(t)\) is known.

\end{englishtranslation}

\end{bilingualpair}

\begin{bilingualpair}

\textbf{分析}：左端是一个卷积，原方程可写为

\nopagebreak[4]\begin{englishtranslation}

\textbf{Analysis.} The left-hand side is a convolution, so the equation
can be written as

\end{englishtranslation}

\[
e^{-|t|}*x(t)=f(t).
\]

\end{bilingualpair}

\begin{bilingualpair}

\textbf{第一步，取傅里叶变换。}

\nopagebreak[4]\begin{englishtranslation}

\textbf{Step 1. Take the Fourier transform.}

\end{englishtranslation}

\end{bilingualpair}

\begin{bilingualpair}

利用卷积定理，

\nopagebreak[4]\begin{englishtranslation}

By the convolution theorem,

\end{englishtranslation}

\[
\mathcal F[e^{-|t|}*x(t)]
=\mathcal F[f(t)],
\]

\end{bilingualpair}

\[
\mathcal F[e^{-|t|}]\cdot\mathcal F[x(t)]
=F(\omega).
\]

\begin{bilingualpair}

即

\nopagebreak[4]\begin{englishtranslation}

That is,

\end{englishtranslation}

\[
\frac2{\omega^2+1}\overline X(\omega)=F(\omega).
\]

\end{bilingualpair}

\begin{bilingualpair}

所以

\nopagebreak[4]\begin{englishtranslation}

Therefore,

\end{englishtranslation}

\[
\overline X(\omega)
=\frac{\omega^2+1}{2}F(\omega).
\]

\end{bilingualpair}

\begin{bilingualpair}

\textbf{第二步，取傅里叶反变换。}

\nopagebreak[4]\begin{englishtranslation}

\textbf{Step 2. Take the inverse Fourier transform.}

\end{englishtranslation}

\[
\begin{aligned}
x(t)
&=\mathcal F^{-1}\!\left[
\frac{\omega^2+1}{2}F(\omega)
\right]\\
&=\frac12\left\{
\mathcal F^{-1}[\omega^2F(\omega)]
+\mathcal F^{-1}[F(\omega)]
\right\}\\
&=\frac12\left\{
-\mathcal F^{-1}[(i\omega)^2F(\omega)]
+f(t)
\right\}\\
&=\frac12[-f''(t)+f(t)].
\end{aligned}
\]

\end{bilingualpair}

\begin{bilingualpair}

\textbf{注：微分性}

\nopagebreak[4]\begin{englishtranslation}

\textbf{Note: the differentiation property}

\end{englishtranslation}

\[
f''(t)\ \xleftrightarrow{\mathcal F}\
(i\omega)^2F(\omega),
\]

\end{bilingualpair}

\begin{bilingualpair}

条件为

\nopagebreak[4]\begin{englishtranslation}

The conditions are

\end{englishtranslation}

\[
f(\pm\infty)=0,
\qquad
f'(\pm\infty)=0.
\]

\end{bilingualpair}

\begin{bilingualpair}

\textbf{特别地}，设

\nopagebreak[4]\begin{englishtranslation}

\textbf{In particular}, let

\end{englishtranslation}

\[
f(t)=e^{-|t|}-\frac12e^{-2|t|},
\]

\end{bilingualpair}

\begin{bilingualpair}

则可以代入上述公式求解。

\nopagebreak[4]\begin{englishtranslation}

The solution can then be obtained by substitution into the formula
above. 
\end{englishtranslation}

\end{bilingualpair}

\begin{bilingualpair}
Notes:
\par
\vspace*{3cm}
\end{bilingualpair}

\Needspace{8\baselineskip}

\subsection{\texorpdfstring{3. 例
7.21：求解正弦核积分方程\\[2pt]{\normalfont\small\color{TranslationBlue} Example
7.21: Solving an Integral Equation with a Sine
Kernel}}{3. 例 7.21：求解正弦核积分方程Example 7.21: Solving an Integral Equation with a Sine Kernel}}

\begin{bilingualpair}

求解

\nopagebreak[4]\begin{englishtranslation}

Solve

\end{englishtranslation}

\[
\int_0^{+\infty}x(\omega)\sin\omega t\,\mathrm d\omega
=e^{-t},
\qquad t>0.
\]

\end{bilingualpair}

\Needspace{8\baselineskip}

\subsubsection{\texorpdfstring{回忆：傅里叶积分公式（反变换）\\[2pt]{\normalfont\small\color{TranslationBlue} Review:
The Fourier Integral Formula (Inverse
Transform)}}{回忆：傅里叶积分公式（反变换）Review: The Fourier Integral Formula (Inverse Transform)}}

\[
f(t)=\frac1{2\pi}
\int_{-\infty}^{+\infty}
\left[
\int_{-\infty}^{+\infty}
f(\tau)e^{-i\omega\tau}\,\mathrm d\tau
\right]e^{i\omega t}\,\mathrm d\omega.
\]

\begin{bilingualpair}

\textbf{特别地}

\nopagebreak[4]\begin{englishtranslation}

\textbf{Special cases}

\end{englishtranslation}

\end{bilingualpair}

\begin{bilingualpair}

\textbf{（I）偶函数}

\nopagebreak[4]\begin{englishtranslation}

\textbf{(I) Even functions}

\end{englishtranslation}

\end{bilingualpair}

\begin{bilingualpair}

若 \(f(-t)=f(t)\)，则

\nopagebreak[4]\begin{englishtranslation}

If \(f(-t)=f(t)\), then

\end{englishtranslation}

\[
f(t)=\frac2\pi
\int_0^{+\infty}
\left[
\int_0^{+\infty}
f(\tau)\cos\omega\tau\,\mathrm d\tau
\right]\cos\omega t\,\mathrm d\omega.
\]

\end{bilingualpair}

\begin{bilingualpair}

称为\textbf{余弦傅里叶积分公式}。

\nopagebreak[4]\begin{englishtranslation}

This is called the \textbf{Fourier cosine integral formula}.

\end{englishtranslation}

\end{bilingualpair}

\begin{bilingualpair}

\textbf{（II）奇函数}

\nopagebreak[4]\begin{englishtranslation}

\textbf{(II) Odd functions}

\end{englishtranslation}

\end{bilingualpair}

\begin{bilingualpair}

若 \(f(-t)=-f(t)\)，则

\nopagebreak[4]\begin{englishtranslation}

If \(f(-t)=-f(t)\), then

\end{englishtranslation}

\[
f(t)=\frac2\pi
\int_0^{+\infty}
\left[
\int_0^{+\infty}
f(\tau)\sin\omega\tau\,\mathrm d\tau
\right]\sin\omega t\,\mathrm d\omega.
\]

\end{bilingualpair}

\begin{bilingualpair}

称为\textbf{正弦傅里叶积分公式}。

\nopagebreak[4]\begin{englishtranslation}

This is called the \textbf{Fourier sine integral formula}.

\end{englishtranslation}

\end{bilingualpair}

\Needspace{8\baselineskip}

\subsubsection{\texorpdfstring{回到原题\\[2pt]{\normalfont\small\color{TranslationBlue} Return
to the Integral Equation}}{回到原题Return to the Integral Equation}}

\begin{bilingualpair}

由

\nopagebreak[4]\begin{englishtranslation}

From

\end{englishtranslation}

\[
\int_0^{+\infty}x(\omega)\sin\omega t\,\mathrm d\omega=e^{-t},
\]

\end{bilingualpair}

\begin{bilingualpair}

得

\nopagebreak[4]\begin{englishtranslation}

we obtain

\end{englishtranslation}

\[
\frac2\pi\int_0^{+\infty}
x(\omega)\sin\omega t\,\mathrm d\omega
=\frac2\pi e^{-t}.
\]

\end{bilingualpair}

\begin{bilingualpair}

由正弦积分公式可知

\nopagebreak[4]\begin{englishtranslation}

The Fourier sine integral formula gives

\end{englishtranslation}

\[
x(\omega)=\int_0^{+\infty}
\frac2\pi e^{-t}\sin\omega t\,\mathrm dt.
\]

\end{bilingualpair}

\begin{bilingualpair}

列出两种计算方法：

\nopagebreak[4]\begin{englishtranslation}

Here suggests two methods of evaluating the integral.

\end{englishtranslation}

\end{bilingualpair}

\begin{itemize}
\item
  \begin{bilingualpair}

  方法一：两次分部积分。

  \nopagebreak[4]\begin{englishtranslation}

  Method 1: integrate by parts twice.

  \end{englishtranslation}

  \end{bilingualpair}
\item
  \begin{bilingualpair}

  方法二：利用 Euler 公式

  \nopagebreak[4]\begin{englishtranslation}

  Method 2: use Euler's formula,

  \end{englishtranslation}

  \end{bilingualpair}
\end{itemize}

\[
\sin\omega t=\frac{e^{i\omega t}-e^{-i\omega t}}{2i}.
\]

\begin{bilingualpair}
Notes:
\par
\vspace*{3cm}
\end{bilingualpair}

\Needspace{12\baselineskip}

\section{\texorpdfstring{三、第 8
章：拉普拉斯变换\\[2pt]{\normalfont\normalsize\color{TranslationBlue} Chapter
8: The Laplace
Transform}}{三、第 8 章：拉普拉斯变换Chapter 8: The Laplace Transform}}

\Needspace{8\baselineskip}

\subsection{\texorpdfstring{1.
拉普拉斯变换的定义\\[2pt]{\normalfont\small\color{TranslationBlue} Definition
of the Laplace
Transform}}{1. 拉普拉斯变换的定义Definition of the Laplace Transform}}

\begin{bilingualpair}

拉普拉斯变换，简称拉氏变换。

\nopagebreak[4]\begin{englishtranslation}

In Chinese, the name ``Laplace transform'' is also abbreviated as
``拉氏变换''.

\end{englishtranslation}

\end{bilingualpair}

\begin{bilingualpair}

对于定义在 \([0,+\infty)\) 上的函数 \(f(t)\)，定义

\nopagebreak[4]\begin{englishtranslation}

For a function \(f(t)\) defined on \([0,+\infty)\), define

\end{englishtranslation}

\[
F(p)=\int_0^{+\infty}f(t)e^{-pt}\,\mathrm dt,
\]

\end{bilingualpair}

\begin{bilingualpair}

其中

\nopagebreak[4]\begin{englishtranslation}

where

\end{englishtranslation}

\[
p=\sigma+i\omega
\]

\end{bilingualpair}

\begin{bilingualpair}

是复变量。

\nopagebreak[4]\begin{englishtranslation}

is a complex variable.

\end{englishtranslation}

\end{bilingualpair}

\begin{bilingualpair}

称 \(\{f(t),F(p)\}\) 为一个拉氏变换对，记作

\nopagebreak[4]\begin{englishtranslation}

The pair \(\{f(t),F(p)\}\) is called a Laplace transform pair. We write

\end{englishtranslation}

\[
\mathcal L[f(t)]=F(p),
\]

\end{bilingualpair}

\begin{bilingualpair}

或

\nopagebreak[4]\begin{englishtranslation}

or

\end{englishtranslation}

\[
\mathcal L^{-1}[F(p)]=f(t).
\]

\end{bilingualpair}

\begin{bilingualpair}
Notes:
\par
\vspace*{3cm}
\end{bilingualpair}

\Needspace{8\baselineskip}

\subsection{\texorpdfstring{2.
指数级函数与拉氏变换存在定理\\[2pt]{\normalfont\small\color{TranslationBlue} Functions
of Exponential Order and the Existence
Theorem}}{2. 指数级函数与拉氏变换存在定理Functions of Exponential Order and the Existence Theorem}}

\Needspace{8\baselineskip}

\subsubsection{\texorpdfstring{定义：指数级函数\\[2pt]{\normalfont\small\color{TranslationBlue} Definition:
A Function of Exponential
Order}}{定义：指数级函数Definition: A Function of Exponential Order}}

\begin{bilingualpair}

若

\nopagebreak[4]\begin{englishtranslation}

If

\end{englishtranslation}

\[
|f(t)|\le Me^{\sigma_ct},
\]

\end{bilingualpair}

\begin{bilingualpair}

则称 \(f(t)\) 是``指数级函数''，\(\sigma_c\) 是其增长指数。
\(\sigma_c\ge0\)。

\nopagebreak[4]\begin{englishtranslation}

then \(f(t)\) is said to be of exponential order, and \(\sigma_c\) is
its growth exponent. \(\sigma_c\ge0\).

\end{englishtranslation}

\end{bilingualpair}

\Needspace{8\baselineskip}

\subsubsection{\texorpdfstring{定理
8.1：拉氏变换存在定理\\[2pt]{\normalfont\small\color{TranslationBlue} Theorem
8.1: Existence of the Laplace
Transform}}{定理 8.1：拉氏变换存在定理Theorem 8.1: Existence of the Laplace Transform}}

\begin{bilingualpair}

若定义在 \([0,+\infty)\) 上的函数 \(f(t)\) 满足

\nopagebreak[4]\begin{englishtranslation}

Suppose that \(f(t)\), defined on \([0,+\infty)\), satisfies the
following conditions.

\end{englishtranslation}

\end{bilingualpair}

\begin{enumerate}
\def\labelenumi{\arabic{enumi}.}
\item
  \begin{bilingualpair}

  \(f(t)\) 满足 Dirichlet 条件；

  \nopagebreak[4]\begin{englishtranslation}

  \(f(t)\) satisfies the Dirichlet conditions.

  \end{englishtranslation}

  \end{bilingualpair}
\item
  \begin{bilingualpair}

  \(f(t)\) 是指数级的，即 \(|f(t)|\le Me^{\sigma_ct}\)；

  \nopagebreak[4]\begin{englishtranslation}

  \(f(t)\) is of exponential order, that is,
  \(|f(t)|\le Me^{\sigma_ct}\).

  \end{englishtranslation}

  \end{bilingualpair}
\end{enumerate}

\begin{bilingualpair}

则

\nopagebreak[4]\begin{englishtranslation}

then

\end{englishtranslation}

\[
F(p)=\int_0^{+\infty}f(t)e^{-pt}\,\mathrm dt
\]

\end{bilingualpair}

\begin{bilingualpair}

在右半平面

\nopagebreak[4]\begin{englishtranslation}

in the right half-plane

\end{englishtranslation}

\[
\operatorname{Re}p>\sigma_c
\]

\end{bilingualpair}

\begin{bilingualpair}

内存在且解析。

\nopagebreak[4]\begin{englishtranslation}

exists and is analytic.

\end{englishtranslation}

\end{bilingualpair}

\begin{bilingualpair}

注：令 \(\widetilde f(t)=f(t)u(t)\)。

\nopagebreak[4]\begin{englishtranslation}

note: let \(\widetilde f(t)=f(t)u(t)\).

\end{englishtranslation}

\end{bilingualpair}

\begin{bilingualpair}

\textbf{原图说明}：复平面以横轴表示
\(\sigma=\operatorname{Re}p\)，纵轴表示
\(\omega=\operatorname{Im}p\)；直线 \(\sigma=\sigma_c\)
右侧的区域用斜线标出，表示收敛与解析的右半平面。

\nopagebreak[4]\begin{englishtranslation}

\textbf{Description of the original figure.} In the complex plane, the
horizontal axis is \(\sigma=\operatorname{Re}p\) and the vertical axis
is \(\omega=\operatorname{Im}p\). The region to the right of
\(\sigma=\sigma_c\) is shaded to indicate the right half-plane of
convergence and analyticity.

\end{englishtranslation}

\end{bilingualpair}

\Needspace{8\baselineskip}

\subsubsection{\texorpdfstring{证明：绝对收敛\\[2pt]{\normalfont\small\color{TranslationBlue} Proof:
Absolute Convergence}}{证明：绝对收敛Proof: Absolute Convergence}}

\begin{bilingualpair}

因为

\nopagebreak[4]\begin{englishtranslation}

Since

\end{englishtranslation}

\[
|f(t)|\le Me^{\sigma_ct}
\]

\end{bilingualpair}

\begin{bilingualpair}

且

\nopagebreak[4]\begin{englishtranslation}

and

\end{englishtranslation}

\[
|e^{-pt}|
=|e^{-\sigma t}e^{-i\omega t}|
=e^{-\sigma t},
\]

\end{bilingualpair}

\begin{bilingualpair}

所以当 \(\sigma>\sigma_c\) 时，

\nopagebreak[4]\begin{englishtranslation}

for \(\sigma>\sigma_c\) we have

\end{englishtranslation}

\[
\begin{aligned}
\int_0^{+\infty}|f(t)e^{-pt}|\,\mathrm dt
&\le\int_0^{+\infty}Me^{\sigma_ct}|e^{-pt}|\,\mathrm dt\\
&=\int_0^{+\infty}Me^{-(\sigma-\sigma_c)t}\,\mathrm dt\\
&=\left.
M\frac{e^{-(\sigma-\sigma_c)t}}{-(\sigma-\sigma_c)}
\right|_0^{+\infty}\\
&=\frac{M}{\sigma-\sigma_c}.
\end{aligned}
\]

\end{bilingualpair}

\begin{bilingualpair}

因此广义积分

\nopagebreak[4]\begin{englishtranslation}

Therefore, the improper integral

\end{englishtranslation}

\[
\int_0^{+\infty}f(t)e^{-pt}\,\mathrm dt
\]

\end{bilingualpair}

\begin{bilingualpair}

存在，且绝对收敛。

\nopagebreak[4]\begin{englishtranslation}

exists and converges absolutely.

\end{englishtranslation}

\end{bilingualpair}

\Needspace{8\baselineskip}

\subsubsection{\texorpdfstring{证明：解析性\\[2pt]{\normalfont\small\color{TranslationBlue} Proof:
Analyticity}}{证明：解析性Proof: Analyticity}}

\begin{bilingualpair}

原稿进一步说明原积分以及

\nopagebreak[4]\begin{englishtranslation}

The original notes further state that both the original integral and

\end{englishtranslation}

\[
\int_0^{+\infty}
f(t)\frac{\partial}{\partial p}(e^{-pt})\,\mathrm dt
\]

\end{bilingualpair}

\begin{bilingualpair}

具有一致收敛性，从而可对参数 \(p\) 求导：

\nopagebreak[4]\begin{englishtranslation}

converge uniformly, allowing differentiation with respect to the
parameter \(p\).

\end{englishtranslation}

\[
\begin{aligned}
F'(p)
&=\frac{\mathrm d}{\mathrm dp}
\int_0^{+\infty}f(t)e^{-pt}\,\mathrm dt\\
&=\int_0^{+\infty}
f(t)\frac{\partial}{\partial p}(e^{-pt})\,\mathrm dt\\
&=\int_0^{+\infty}(-t)f(t)e^{-pt}\,\mathrm dt.
\end{aligned}
\]

\end{bilingualpair}

\begin{bilingualpair}

故 \(F(p)\) 在 \(\operatorname{Re}p>\sigma_c\) 内解析。

\nopagebreak[4]\begin{englishtranslation}

Hence \(F(p)\) is analytic in \(\operatorname{Re}p>\sigma_c\).

\end{englishtranslation}

\end{bilingualpair}

\begin{bilingualpair}
Notes:
\par
\vspace*{3cm}
\end{bilingualpair}

\Needspace{8\baselineskip}

\subsubsection{\texorpdfstring{原稿补充说明\\[2pt]{\normalfont\small\color{TranslationBlue} Additional
Remarks in the Original
Notes}}{原稿补充说明Additional Remarks in the Original Notes}}

\begin{enumerate}
\def\labelenumi{\arabic{enumi}.}
\item
  \begin{bilingualpair}

  ``由上述证明中知，\(F(p)\) 解析且 \(F(+\infty)=0\)。''

  \nopagebreak[4]\begin{englishtranslation}

  ``The preceding proof shows that \(F(p)\) is analytic and
  \(F(+\infty)=0\).''

  \end{englishtranslation}

  \end{bilingualpair}
\item
  \begin{bilingualpair}

  ``由解析延拓，\(F(p)\) 于 \(p\) 平面上除若干孤立奇点外都解析。''

  \nopagebreak[4]\begin{englishtranslation}

  ``By analytic continuation, \(F(p)\) is analytic throughout the
  \(p\)-plane except at certain isolated singularities.''

  \end{englishtranslation}

  \end{bilingualpair}
\end{enumerate}

\begin{bilingualpair}
Notes:
\par
\vspace*{3cm}
\end{bilingualpair}

\Needspace{8\baselineskip}

\subsection{\texorpdfstring{3.
常用拉氏变换对\\[2pt]{\normalfont\small\color{TranslationBlue} Common
Laplace Transform
Pairs}}{3. 常用拉氏变换对Common Laplace Transform Pairs}}

\begin{bilingualpair}

\emph{教材 P.345，共 73对''。}

\nopagebreak[4]\begin{englishtranslation}

\emph{textbook, p.~345.``73 pairs in total''.}

\end{englishtranslation}

\end{bilingualpair}

\Needspace{8\baselineskip}

\subsubsection{\texorpdfstring{第 1
对：单位阶跃函数\\[2pt]{\normalfont\small\color{TranslationBlue} Pair 1:
The Unit Step
Function}}{第 1 对：单位阶跃函数Pair 1: The Unit Step Function}}

\[
u(t)\ \xleftrightarrow{\mathcal L}\ \frac1p,
\qquad \operatorname{Re}p>0.
\]

\begin{bilingualpair}

因为

\nopagebreak[4]\begin{englishtranslation}

Since

\end{englishtranslation}

\[
\begin{aligned}
\mathcal L[u(t)]
&=\int_0^{+\infty}1\cdot e^{-pt}\,\mathrm dt\\
&=\left.\frac{e^{-pt}}{-p}\right|_0^{+\infty}\\
&=\frac1p.
\end{aligned}
\]

\end{bilingualpair}

\Needspace{8\baselineskip}

\subsubsection{\texorpdfstring{第 2
对：幂函数\\[2pt]{\normalfont\small\color{TranslationBlue} Pair 2: Power
Functions}}{第 2 对：幂函数Pair 2: Power Functions}}

\begin{bilingualpair}

类似地，

\nopagebreak[4]\begin{englishtranslation}

Similarly,

\end{englishtranslation}

\[
\begin{aligned}
\mathcal L[t]
&=\int_0^{+\infty}te^{-pt}\,\mathrm dt\\
&=\int_0^{+\infty}t\,\mathrm d\!\left(\frac{e^{-pt}}{-p}\right)\\
&=\left.-\frac{te^{-pt}}p\right|_0^{+\infty}
+\int_0^{+\infty}\frac{e^{-pt}}p\,\mathrm dt\\
&=0+\left.\frac{e^{-pt}}{-p^2}\right|_0^{+\infty}\\
&=\frac1{p^2}.
\end{aligned}
\]

\end{bilingualpair}

\[
\begin{aligned}
\mathcal L[t^2]
&=\int_0^{+\infty}t^2e^{-pt}\,\mathrm dt\\
&=\left.t^2\frac{e^{-pt}}{-p}\right|_0^{+\infty}
-\int_0^{+\infty}\frac{e^{-pt}}{-p}\,2t\,\mathrm dt\\
&=\frac2p\cdot\frac1{p^2}\\
&=\frac{2!}{p^3}.
\end{aligned}
\]

\begin{bilingualpair}

推广为

\nopagebreak[4]\begin{englishtranslation}

Generalizing gives

\end{englishtranslation}

\[
\mathcal L[t^n]=\frac{n!}{p^{n+1}},
\]

\end{bilingualpair}

\begin{bilingualpair}

即

\nopagebreak[4]\begin{englishtranslation}

That is,

\end{englishtranslation}

\[
t^n\ \xleftrightarrow{\mathcal L}\
\frac{n!}{p^{n+1}},
\qquad n=0,1,2,\ldots,\quad \operatorname{Re}p>0.
\]

\end{bilingualpair}

\Needspace{8\baselineskip}

\subsubsection{\texorpdfstring{第 4
对：指数函数\\[2pt]{\normalfont\small\color{TranslationBlue} Pair 4: An
Exponential
Function}}{第 4 对：指数函数Pair 4: An Exponential Function}}

\[
e^{at}\ \xleftrightarrow{\mathcal L}\ \frac1{p-a},
\]

\begin{bilingualpair}

其中 \(a\) 为常数。

\nopagebreak[4]\begin{englishtranslation}

where \(a\) is a constant.

\end{englishtranslation}

\[
\begin{aligned}
\mathcal L[e^{at}]
&=\int_0^{+\infty}e^{at}e^{-pt}\,\mathrm dt\\
&=\left.\frac{e^{(a-p)t}}{a-p}\right|_0^{+\infty}\\
&=\frac1{p-a},
\end{aligned}
\]

\end{bilingualpair}

\begin{bilingualpair}

收敛条件为

\nopagebreak[4]\begin{englishtranslation}

The convergence condition is

\end{englishtranslation}

\[
\operatorname{Re}(a-p)<0
\quad\Longleftrightarrow\quad
\operatorname{Re}p>\operatorname{Re}a.
\]

\end{bilingualpair}

\Needspace{8\baselineskip}

\subsubsection{\texorpdfstring{第 20---23
对：三角函数与双曲函数\\[2pt]{\normalfont\small\color{TranslationBlue} Pairs
20--23: Trigonometric and Hyperbolic
Functions}}{第 20---23 对：三角函数与双曲函数Pairs 20--23: Trigonometric and Hyperbolic Functions}}

\Needspace{20\baselineskip}\begingroup\small\renewcommand{\arraystretch}{1.65}\everymath{\displaystyle}

\begin{longtable}[]{@{}
  >{\raggedright\arraybackslash}p{(\columnwidth - 4\tabcolsep) * \real{0.1600}}
  >{\raggedright\arraybackslash}p{(\columnwidth - 4\tabcolsep) * \real{0.2800}}
  >{\raggedright\arraybackslash}p{(\columnwidth - 4\tabcolsep) * \real{0.5600}}@{}}
\toprule\noalign{}
\begin{minipage}[b]{\linewidth}\raggedright
原编号 / Original no.
\end{minipage} & \begin{minipage}[b]{\linewidth}\raggedright
原函数 / Function
\end{minipage} & \begin{minipage}[b]{\linewidth}\raggedright
拉氏变换 / Laplace transform
\end{minipage} \\
\midrule\noalign{}
\endhead
\bottomrule\noalign{}
\endlastfoot
20 & \(\sin\beta t\) & \(\dfrac{\beta}{p^2+\beta^2}\) \\
21 & \(\cos\beta t\) & \(\dfrac{p}{p^2+\beta^2}\) \\
22 & \(\sinh\beta t\) & \(\dfrac{\beta}{p^2-\beta^2}\) \\
23 & \(\cosh\beta t\) & \(\dfrac{p}{p^2-\beta^2}\) \\
\end{longtable}

\endgroup

\begin{editorialnote}

\begin{bilingualpair}

\textbf{注}：当 \(\beta\)为实数时，前两式可取 \(\operatorname{Re}p>0\)，后两式可取
\(\operatorname{Re}p>|\beta|\)。

\nopagebreak[4]\begin{englishtranslation}

\textbf{Editorial note.} For real \(\beta\), the first two formulas hold for
\(\operatorname{Re}p>0\), and the last two hold for \(\operatorname{Re}p>|\beta|\).

\end{englishtranslation}

\end{bilingualpair}

\end{editorialnote}

\begin{bilingualpair}

\textbf{以正弦函数为例}

\nopagebreak[4]\begin{englishtranslation}

\textbf{Derivation for the sine function}

\end{englishtranslation}

\[
\begin{aligned}
\mathcal L[\sin\beta t]
&=\mathcal L\!\left[
\frac{e^{i\beta t}-e^{-i\beta t}}{2i}
\right]\\
&=\frac1{2i}\left\{
\mathcal L[e^{i\beta t}]
-\mathcal L[e^{-i\beta t}]
\right\}\\
&=\frac1{2i}\left[
\frac1{p-i\beta}-\frac1{p+i\beta}
\right]\\
&=\frac1{2i}\cdot\frac{2i\beta}{p^2+\beta^2}\\
&=\frac{\beta}{p^2+\beta^2}.
\end{aligned}
\]

\end{bilingualpair}

\begin{bilingualpair}
Notes:
\par
\vspace*{3cm}
\end{bilingualpair}

\Needspace{12\baselineskip}

\section{\texorpdfstring{四、拉氏变换的性质\\[2pt]{\normalfont\normalsize\color{TranslationBlue} Properties
of the Laplace
Transform}}{四、拉氏变换的性质Properties of the Laplace Transform}}

\begin{bilingualpair}

\emph{教材 P.214。}

\nopagebreak[4]\begin{englishtranslation}

\emph{textbook, p.~214.}

\end{englishtranslation}

\end{bilingualpair}

\Needspace{8\baselineskip}

\subsection{\texorpdfstring{1. 性质
1：线性性\\[2pt]{\normalfont\small\color{TranslationBlue} Property 1:
Linearity}}{1. 性质 1：线性性Property 1: Linearity}}

\[
\mathcal L[af_1+bf_2]
=a\mathcal L[f_1]+b\mathcal L[f_2].
\]

\[
\mathcal L^{-1}[aF_1+bF_2]
=a\mathcal L^{-1}[F_1]+b\mathcal L^{-1}[F_2].
\]

\Needspace{8\baselineskip}

\subsubsection{\texorpdfstring{例：求拉氏反变换\\[2pt]{\normalfont\small\color{TranslationBlue} Example:
Find the Inverse Laplace
Transform}}{例：求拉氏反变换Example: Find the Inverse Laplace Transform}}

\[
\mathcal L^{-1}\!\left[
\frac{p}{(p-1)(p^2+4)}
\right].
\]

\begin{bilingualpair}

先作部分分式分解：

\nopagebreak[4]\begin{englishtranslation}

First, decompose the expression into partial fractions.

\end{englishtranslation}

\[
\begin{aligned}
\frac{p}{(p-1)(p^2+4)}
&=\frac{A}{p-1}+\frac{Bp+C}{p^2+4}\\
&=\frac{1/5}{p-1}
+\frac{-p/5+4/5}{p^2+4}.
\end{aligned}
\]

\end{bilingualpair}

\begin{bilingualpair}

所以

\nopagebreak[4]\begin{englishtranslation}

Therefore,

\end{englishtranslation}

\[
\begin{aligned}
\mathcal L^{-1}\!\left[
\frac{p}{(p-1)(p^2+4)}
\right]
&=\frac15\mathcal L^{-1}\!\left[\frac1{p-1}\right]
-\frac15\mathcal L^{-1}\!\left[\frac{p}{p^2+4}\right]\\
&\quad+\frac45\mathcal L^{-1}\!\left[\frac1{p^2+4}\right]\\
&=\frac15e^t-\frac15\cos2t+\frac25\sin2t.
\end{aligned}
\]

\end{bilingualpair}

\begin{bilingualpair}

原稿红字提示所用变换对：

\nopagebreak[4]\begin{englishtranslation}

The original notes highlight the following transform pairs in red.

\end{englishtranslation}

\[
e^{at}\ \xleftrightarrow{\mathcal L}\ \frac1{p-a},
\]

\end{bilingualpair}

\[
\cos at\ \xleftrightarrow{\mathcal L}\ \frac{p}{p^2+a^2},
\]

\[
\sin at\ \xleftrightarrow{\mathcal L}\ \frac{a}{p^2+a^2}.
\]

\begin{bilingualpair}
Notes:
\par
\vspace*{3cm}
\end{bilingualpair}

\Needspace{8\baselineskip}

\subsection{\texorpdfstring{2. 性质
2：相似性\\[2pt]{\normalfont\small\color{TranslationBlue} Property 2:
Scaling}}{2. 性质 2：相似性Property 2: Scaling}}

\begin{bilingualpair}

设

\nopagebreak[4]\begin{englishtranslation}

Let

\end{englishtranslation}

\[
f(t)\ \xleftrightarrow{\mathcal L}\ F(p).
\]

\end{bilingualpair}

\begin{bilingualpair}

则

\nopagebreak[4]\begin{englishtranslation}

then

\end{englishtranslation}

\[
f(at)\ \xleftrightarrow{\mathcal L}\
\frac1aF\!\left(\frac pa\right),
\qquad a>0.
\]

\end{bilingualpair}

\begin{bilingualpair}

以及

\nopagebreak[4]\begin{englishtranslation}

and

\end{englishtranslation}

\[
\frac1a f\!\left(\frac ta\right)
\ \xleftrightarrow{\mathcal L}\
F(ap).
\]

\end{bilingualpair}

\begin{bilingualpair}
Notes:
\par
\vspace*{3cm}
\end{bilingualpair}

\Needspace{8\baselineskip}

\subsection{\texorpdfstring{3. 性质
3：延迟性\\[2pt]{\normalfont\small\color{TranslationBlue} Property 3:
Time Delay}}{3. 性质 3：延迟性Property 3: Time Delay}}

\begin{bilingualpair}

\emph{``对 \(f(t)u(t)\) 的平移性''。}

\nopagebreak[4]\begin{englishtranslation}

\emph{``shifting \(f(t)u(t)\)''.}

\end{englishtranslation}

\[
f(t-t_0)u(t-t_0)
\ \xleftrightarrow{\mathcal L}\
e^{-pt_0}F(p).
\]

\end{bilingualpair}

\begin{bilingualpair}

即

\nopagebreak[4]\begin{englishtranslation}

That is,

\end{englishtranslation}

\[
\mathcal L[f(t-t_0)u(t-t_0)]
=e^{-pt_0}F(p),
\]

\end{bilingualpair}

\begin{bilingualpair}

以及

\nopagebreak[4]\begin{englishtranslation}

and

\end{englishtranslation}

\[
\mathcal L^{-1}[e^{-pt_0}F(p)]
=f(t-t_0)u(t-t_0).
\]

\end{bilingualpair}

\begin{editorialnote}

\begin{bilingualpair}

\textbf{注}：单边拉普拉斯变换的这一延迟公式通常取
\(t_0\ge0\)。

\nopagebreak[4]\begin{englishtranslation}

\textbf{note.} For the unilateral Laplace transform, this
time-delay formula is normally used with \(t_0\ge0\).

\end{englishtranslation}

\end{bilingualpair}

\end{editorialnote}

\begin{bilingualpair}
Notes:
\par
\vspace*{3cm}
\end{bilingualpair}

\Needspace{8\baselineskip}

\subsection{\texorpdfstring{4. 性质 4：对 \(F(p)\)
的平移性\\[2pt]{\normalfont\small\color{TranslationBlue} Property 4:
Shifting in the Transform
Variable}}{4. 性质 4：对 F(p) 的平移性Property 4: Shifting in the Transform Variable}}

\begin{bilingualpair}

\[
e^{p_0t}f(t)
\ \xleftrightarrow{\mathcal L}\
F(p-p_0).
\]

\end{bilingualpair}

\begin{bilingualpair}

原稿证明从定义开始：

\nopagebreak[4]\begin{englishtranslation}

The original proof starts from the definition.

\end{englishtranslation}

\[
F(p)=\int_0^{+\infty}f(t)e^{-pt}\,\mathrm dt.
\]

\end{bilingualpair}

\begin{bilingualpair}

本页到此未继续展开证明。

\end{bilingualpair}

\begin{bilingualpair}
Notes:
\par
\vspace*{3cm}
\end{bilingualpair}

\Needspace{12\baselineskip}

\section{\texorpdfstring{五、12.1
作业\\[2pt]{\normalfont\normalsize\color{TranslationBlue} Assignment for
December 1}}{五、12.1 作业Assignment for December 1}}

\begin{itemize}
\item
  \begin{bilingualpair}

  教材 P.204，A 组第 9、10 题。

  \nopagebreak[4]\begin{englishtranslation}

  Textbook, p.~204: Group A, Exercises 9 and 10.

  \end{englishtranslation}

  \end{bilingualpair}
\item
  \begin{bilingualpair}

  教材 P.204，B 组第 1、2、5、6、8 题。

  \nopagebreak[4]\begin{englishtranslation}

  Textbook, p.~204: Group B, Exercises 1, 2, 5, 6, and 8.

  \end{englishtranslation}

  \end{bilingualpair}
\item
  \begin{bilingualpair}

  教材 P.236，第 1 题。

  \nopagebreak[4]\begin{englishtranslation}

  Textbook, p.~236: Exercise 1.

  \end{englishtranslation}

  \end{bilingualpair}
\end{itemize}

\begin{bilingualpair}
Notes:
\end{bilingualpair}

\end{document}
